In questo capitolo si presenta il consuntivo delle attività e delle ore svolte, confrontato con il preventivo iniziale (Tabella~\ref{tab:pianificazione-attivita}). Gli scostamenti sono principalmente imputabili ai limiti informativi di \gls{adrr} sui settori elementari, che hanno richiesto un maggiore impegno nell’implementazione delle procedure di esportazione; la ridotta copertura di casi reali ha inoltre comportato una lieve riduzione dei test di performance dei modelli \gls{dac}. In luogo dell’interfaccia utente, è stata privilegiata la generazione di istanze fittizie, riadattando le \gls{fir} per simulare gruppi settoriali e consentire benchmark consistenti. Il totale consuntivato ammonta a 327 ore e i dettagli sono riportati in Tabella~\ref{tab:consuntivo-attivita}.
Di seguito si riporta lo stato degli obiettivi definiti in Capitolo \ref{ch:obiettivi} (cfr. Tabella~\ref{tab:consuntivo-obiettivi}).  

\begin{itemize}
    \item \textit{O01 – O04}: Tutti gli obiettivi obbligatori sono stati raggiunti come previsto
    \item \textit{D01}: Parzialmente raggiunto a causa della copertura limitata di \gls{adrr} sui settori elementari, che ha impedito di realizzare un benchmark completamente basato su dati reali
    \item \textit{D02}: Raggiunto, sebbene con un numero ridotto di test rispetto al preventivo, per via della limitata disponibilità di istanze reali
    \item \textit{F01}: Non raggiunto; al suo posto è stata implementata la generazione di istanze fittizie per consentire il completamento dei benchmark.
\end{itemize}
\clearpage
    \begin{table}[htbp]
\begin{tabularx}{\textwidth}{!{\vrule}c!{\vrule}X!{\vrule}c!{\vrule}}
\hline
\textbf{Obiettivo} & \textbf{Definizione} & \textbf{Stato} \\\hline
\textit{O01} & Acquisizione di competenze su dati concernenti il traffico aereo e il suo controllo & Raggiunto \\
\textit{O02} & Acquisizione dei dati \acrshort{adrr} & Raggiunto \\
\textit{O03} & Implementazione di procedure di esportazione dei dati \acrshort{adrr} nei formati richiesti dai modelli di ottimizzazione per \acrshort{dac} & Raggiunto \\
\textit{O04} & Redazione della documentazione di un manuale utente per le procedure realizzate & Raggiunto \\
\hline
\textit{D01} & Realizzazione di un benchmark di istanze reali DAC & Parzialmente raggiunto \\
\textit{D02} & Test di performance dei modelli di ottimizzazione per \acrshort{dac} sul benchmark realizzato & Raggiunto \\
\hline
\textit{F01} & Implementazione di un'interfaccia utente per l'esecuzione delle procedure di importazione ed esportazione dei dati \acrshort{adrr} & Non raggiunto \\
\hline
\end{tabularx}
\caption{Consuntivo degli obiettivi.}
\label{tab:consuntivo-obiettivi}
\end{table}
   
\begin{table}[H]
\begin{tabularx}{\textwidth}{!{\vrule}c!{\vrule}X!{\vrule}}
	\hline
	\textbf{Durata in ore} & \textbf{Descrizione dell'attività} \\\hline
	
	\textbf{40} & \textbf{Formazione} \\ \hdashline 
    \multirow{3}{0cm}\\ 
    \textit{16} & 
    \textit{Formazione sul traffico e il controllo aereo} \\
    \textit{16} & 
    \textit{Formazione sui modelli per \gls{dac} e sulle tecnologie di ottimizzazione} \\
    \textit{8} & 
    \textit{Formazione su data repository \citet{eurocontrol2025}} \\
    \hline
    
    \textbf{72} & \textbf{Acquisizione dei dati \gls{adrr}}  \\ \hdashline 
    \multirow{4}{0cm}\\ 
    \textit{8} & 
    \textit{Accesso a \gls{adrr} e download dei dati e della documentazione} \\
    \textit{24} & 
    \textit{Studio della documentazione e definizione dei dati per \gls{dac}} \\
    \textit{24} & 
    \textit{Implementazione prototipale di procedure di lettura, analisi e esportazione di dati \gls{adrr}} \\
    \textit{16} & 
    \textit{Redazione documenti} \\
    \hline

    \textbf{40} & \textbf{Progettazione delle procedure di esportazione}  \\ \hdashline 
    \multirow{4}{0cm}\\ 
    \textit{20} & 
    \textit{Definizione dati richiesti da \gls{dac} e di eventuali pre-elaborazioni dei dati \gls{adrr}} \\
    \textit{20} & 
    \textit{Progetto dei moduli di esportazione} \\
    \hline
 
    \textbf{100} (+12) & \textbf{Implementazione delle procedure di esportazione}  \\ \hdashline 
    \multirow{4}{0cm}\\ 
    \textit{84} (+12) & 
    \textit{Implementazione dei moduli di pre-elaborazione e di esportazione} \\
    \textit{16} & 
    \textit{Test e validazione delle procedure} \\
    \hline   

    \textbf{35} (-5) & \textbf{Test di performance dei modelli \gls{dac}}  \\ \hdashline 
    \multirow{4}{0cm}\\ 
    \textit{15} (-5) & 
    \textit{Realizzazione di un benchmark di istanze \gls{dac}} \\
    \textit{20} & 
    \textit{Esecuzione dei test sui modelli di ottimizzazione per \gls{dac}} \\
    \hline   

	\textbf{20} & \textbf{Generazione di istanze fittizie (per benchmark \gls{dac})}  \\ 
    \hline   

    \textbf{20} & \textbf{Redazione documenti conclusivi}  \\ 
    \hline   
    
    \hline \hline
    \textbf{\textit{327 (+7)}} & \multicolumn{1}{|l|}{\textbf{\textit{Totale ore}}} \\\hline
\end{tabularx}
\caption{Consuntivo delle ore di lavoro per attività. Tra parentesi lo scostamento rispetto al preventivo (Tabella~\ref{tab:pianificazione-attivita}).}
\label{tab:consuntivo-attivita}
\end{table}



\section{Considerazioni personali sull’esperienza}
\subsection{Difficoltà incontrate}
Le principali difficoltà hanno riguardato l’allineamento tra ciò che sarebbe “ideale” per il \gls{dac} e ciò che è realmente osservabile in \gls{adrr}. L’assenza di settori elementari e di confini \gls{acc} ha imposto scelte di modellazione conservative e un maggiore dispendio di tempo nell’analisi sistematica della repository alla ricerca di proxy informativi. Nonostante il lavoro di messa a punto delle metriche temporali (\emph{throughput}, \emph{occupancy}) e l’esplorazione approfondita del dataset non abbia prodotto evidenze sufficienti a descrivere o inferire i settori elementari, le procedure sviluppate mantengono una forte utilità metodologica: esse permettono di costruire istanze reali di \gls{dac} non appena i dati necessari diventino disponibili, garantendo la replicabilità e l’adattabilità dei benchmark.

\subsection{Retrospettiva e lezioni apprese}
Lo stage mi ha imposto di integrare ingegneria del dato e modellazione: una definizione elegante è sterile se i dati non la sostengono; simmetricamente, una pipeline robusta deve esplicitare le assunzioni che introduce. Ho constatato che un modello convincente è inutile se le variabili chiave non sono osservabili (nel nostro caso i settori elementari in \gls{adrr}) e che una pipeline apparentemente solida può fuorviare se non rende chiare e tracciabili le scelte operative (ampiezza delle finestre temporali, regole di conteggio).
Non è stato soltanto esercizio di programmazione: mi ha portato oltre il perimetro tipico del corso di laurea, chiedendomi di coniugare metodo e analisi dei dati. Ho dovuto rendere esplicite le ipotesi, documentarle nel codice e verificarle sistematicamente sul dataset. Ho imparato a esplicitare precocemente i limiti osservativi (ciò che non si può misurare), così da distinguere con rapidità gli errori d’implementazione dalle opzioni modellistiche.

